\scanpage{100}
\mspage{\deleted{\unsure{43}}44}
% archive sheet number: 99
Notons la variante bien naturelle

\uline{L 6} \uline{bis} : Comme L6, \deleted{\unsure{avec}} mais en supposant seulement $Y'\mathbin{\unsure{-|\to}}Y$ une $\mathrm{Hot}$-fibration.

Il y a des chances que \inserted{pratiquement,} quand L6 (ou seulement L5) est vérifié, L6 bis l'est aussi.

On va introduire un axiome

\uline{L 7} a) \unsure{Soient|Soit} $(f_i:X_i\mathbin{\unsure{-|\to}}Y_i)_{i\in I}$ une famille de flèches de $\mathrm{Cat}$, soit $f:\coprod X_i\longrightarrow\coprod Y_i$ la flèche \deleted{\illegible{1}} corr. sur les catégories sommes. Pour que $f\in W$, \deleted{\unsure{on}} il f. et s que $f_i\in W$ pour tt $i\in I$.

b) Soit $f:X\longrightarrow Y$ dans $W$. Si $Y$ est $0$-connexe, $X$ aussi.

Cet axiome équivaut à la conjonction de \unsure{L o)|L a)} ($f\in W\Longrightarrow\pi_0(f)$ bijectif) et de la validité du ``lemme'' page 23, que je \unsure{remets|recopie|récris} ici pour mémoire :

\uline{Lemme} L'axiome L7 précédent sur le localiseur fondamental $W$ équivaut : la \deleted{conjonction des} \inserted{conditions suivantes :} $f:X\longrightarrow Y$ \deleted{conditions suivantes} est dans $W$ ssi $\pi_0(f)$ est bijectif, et \deleted{\unsure{a')|ii)}} $f$ induit des $W$-équivalences $f_i:X_i\mathbin{\unsure{-|\to}}Y_i$ sur les composantes connexes.

Il me faut revenir encore sur le 1) de la proposition, il n'est pas vrai que \deleted{se réduise à} \inserted{``c'est''} la condition L4 \deleted{sur} $W$ (que j'admettais, ai-je dit), c'est L4 dans le cas particulier où $Y=S$. Réflexion faite, je la vois comme une variante renforcée de L4, qui ne semble pas \unsure{devoir} découler de L4 ; \deleted{\unsure{Soit quoi} C'est une variante affaiblie de \unsure{l'én.}}

\uline{L 4 bis} Soit $f$ :

\diagram{\begin{tikzcd}[ampersand replacement=\&,column sep=small,row sep=tiny]
 X\arrow[rr]\arrow[dr] \& \& Y\arrow[dl]\\
 \& S\&
\end{tikzcd}}

, \deleted{\unsure{suppos}} dans $\mathrm{Cat}$, supposons que \deleted{\unsure{soit}} $f$ soit une $W$-équiv. \uline{loc. sur} $S$ (i.e. $f_{/s}:X_{/s}\mathbin{\unsure{-|\to}}Y_{/s}$ \unsure{une} $W$-équiv. $\forall s\in S$), alors $f$ est une $W$-équivalence.

% slanted
\margin{équivaut à l'axiome dual, \deleted{\unsure{en prenant}} \unsure{au|ou} \unsure{(cf.)} \deleted{\unsure{lieu de}} \unsure{que} \deleted{\unsure{cela signifie que les foncteurs|en prenant les catégories opposées}}}

